\section{Equilibrio uniforme con marcas de masa variable}
Sea $z=jD_0$, con $D_0$ el escalado bistocástico de $D$. En su soporte $z$ factoriza
en dos escalas positivas. $z$ es acotado y $D_0$ tiene brecha por R2 y la herencia de R7.
En cada eje la medida basal $\mu^0$ asigna masa $s$ a cada índice. Para una marca del
grupo $g=1,2$ se desdobla su índice en un tipo ordinario de masa $s-h_g$ y una copia
de masa $h_g$; en ambos se hereda exactamente la misma fila o columna del núcleo $z$.
Sólo la copia tiene $t=1-x_i$. Denotamos por $\mu$ la medida desdoblada.
En los dos lados la masa de $t\mu$ es idéntica para todo $x$.
El caso físico es $h_1=h_2=s$. Permitimos $|h_g|,|\sigma|\le c$ y $|x_i|\le2$
con $c$ pequeño independiente de $j$, y resolvemos
\[
 v(1-\sigma t_r t_c v)=z\exp(\alpha_r+\beta_c),\qquad
 \int t_c v_{rc}\,d\mu_C=1,\qquad \int t_r v_{rc}\,d\mu_R=1 .
\]
La raíz pequeña en $\sigma$ puede escribirse $v=zq$, con
$q=\exp(\alpha_r+\beta_c)/(1-\sigma t_rt_c v)$ también definido en ceros.
En $(\sigma,h)=(0,0)$ los potenciales son cero, $v=z$.
El jacobiano de márgenes en los tipos ordinarios tiene identidad en sus bloques diagonales
y $D_0,D_0^t$ en los cruzados. En el complemento del modo nulo la brecha da inverso
acotado en $L^2(\mu^0)$. La ecuación de cada coordenada y la cota de $z$ dan la cota en
supremo del inverso para segundo miembro acotado. En la gauge de medias basales iguales
se añade una incógnita escalar con signos opuestos en las ecuaciones de los dos lados.
La compatibilidad por integración contra $t\mu$ fuerza ese escalar a cero; la masa de
cada lado es $1+O(c)$. Los tipos de masa cero tienen ecuación con diagonal uno, cuyo
segundo miembro depende de los tipos ordinarios. No agregan una dirección nula.
La perturbación del mapa y de su jacobiano en bolas pequeñas es $O(|\sigma|+|h_1|+|h_2|)$
en supremo, por variación total de las medidas y la fórmula escalar de la raíz.
Newton con el inverso inicial congelado en una bola fija menor da la solución holomorfa
uniforme; sus normas y $q,1/q$ están acotados. Las ramas logarítmicas se fijan en el origen.
En los parámetros físicos se obtiene la misma rama estacionaria que en el perfil exterior;
en márgenes reales positivos pequeños por casilla se conserva el máximo (concavidad por fila
de torres para suma de fila a lo sumo uno, como en R9). En el límite de margen cero se
descarta la fila o columna contraída. El desdoblamiento se utiliza sólo como prueba analítica.
Escribimos $M_a=t^a\mu$ en cada eje. Cuando un integrando es igual sobre las copias,
la medida proyectada es $\mu^0+\sum_i h_{g(i)}[(1-x_i)^a-1]\delta_i$.
En $\sigma=0$ la solución $z_x=v(0,h,x)$ es igual sobre las copias y sólo depende de $M_1$.
Depende por tanto holomórficamente de las cuatro cantidades $h_{g(i)}x_i$.
